% Section 12.2, from lectures/L12/README.md: the course's learning outcomes, the questions to test
% yourself with, and, as this is the last chapter, "Efter kursen" in the place of the next lecture.

\section{Sammanfattning}
\label{c12:sec:summary}

\needspace{6\baselineskip}
\subsection*{Det här ska du kunna nu}
Kursplanens lärandemål. Efter kursen ska du kunna:
\begin{itemize}
  \item förklara grindar och logiska funktioner;
  \item beskriva binära, hexadecimala och decimala talsystem, och omvandla talvärden mellan dem;
  \item beskriva logisk algebra och förenklingar av booleska uttryck;
  \item redogöra för mikrodatorns uppbyggnad på blockschemanivå;
  \item syntetisera digitala kombinatoriska nät med IC-kretsar och kontakter;
  \item använda Karnaughdiagram för analys och minimering av booleska uttryck;
  \item använda assemblerspråk för att programmera en mikrodator, och programmera den att styra
    styrobjekt;
  \item analysera digitala kombinatoriska nät.
\end{itemize}

\needspace{6\baselineskip}
\subsection*{Frågor att testa dig själv med}
\begin{enumerate}
  \item Vilket tal är \code{0xB6} decimalt, och hur ser det ut binärt?
  \item Vad säger De Morgans lagar, och varför kan varje logisk funktion byggas av enbart
    NAND-grindar?
  \item Varför står raderna i ett Karnaughdiagram i ordningen 00, 01, 11, 10?
  \item Vad gör programräknaren, och vad händer med den vid ett \code{rcall}?
  \item Efter \code{ldi r16, 200} och \code{subi r16, 201}: vad är \code{r16}, och vad är
    \code{C}?
  \item En loop med \code{dec} och \code{brne} går 100 varv. Hur många cykler tar den, och hur lång
    tid är det vid 16~MHz?
  \item Varför läses en nedtryckt knapp som noll, och vad behövs för att den ska läsas som något
    alls när den är släppt?
\end{enumerate}

\needspace{6\baselineskip}
\subsection*{Efter kursen}
Kursen har gått från ettor och nollor till en mikrodator som styr ett trafikljus. Varje steg på
vägen finns kvar i det du har byggt: grindarna finns i ALU:n, vipporna i registren och
programräknaren, kontaktnätens logik i programmets villkorliga hopp.

Vill du vidare finns några naturliga nästa steg:
\begin{itemize}
  \item \textbf{PLC-programmering.} Kontaktnät och stegdiagram från kapitel~\ref{c2:ch}, med samma
    idé om ingångar, beslut och utgångar som i kapitel~\ref{c11:ch}, i den form industrin
    använder.
  \item \textbf{C för mikrokontroller.} Samma ATmega328P och samma register, \code{DDRB},
    \code{PORTB} och \code{PIND}, men i ett språk där kompilatorn skriver assemblern åt dig. Den
    som har skrivit assembler vet vad kompilatorn gör.
  \item \textbf{Programmerbar logik.} Grindnät och vippor beskrivna i ett hårdvarubeskrivande
    språk, och körda på en FPGA.
\end{itemize}

Spara projekten från Labb~3. Ett program som fungerar är den bästa referensen när nästa program
ska skrivas.
